% DIN 8580: the six main groups of manufacturing processes, and where AM sits
\documentclass[tikz,border=6pt]{standalone}
\usepackage{xfab-tikz}
\begin{document}
\begin{tikzpicture}[x=1cm,y=1cm]
\def\w{2.4}\def\g{0.25}
% column x positions
\foreach \k in {0,...,5}{ \pgfmathsetmacro{\x}{\k*(\w+\g)} \coordinate (c\k) at (\x,0); }

% ---------- row 1: what happens to the shape
\node[box=steel, minimum width=\w cm, minimum height=0.7cm, font=\bfseries\small]
  at ($(c0)+(\w/2,3.7)$) {Shape created};
\node[box=steel, minimum width=4*\w cm+3*\g cm, minimum height=0.7cm, font=\bfseries\small]
  at ($(c1)!0.5!(c4)+(\w/2,3.7)$) {Shape changed};
\node[box=steel, minimum width=\w cm, minimum height=0.7cm, font=\bfseries\small]
  at ($(c5)+(\w/2,3.7)$) {Material changed};
% ---------- row 2: cohesion
\foreach \k/\t in {0/created,1/kept,2/reduced,5/{\\at atomic scale}}{
  \node[box=steellight, minimum width=\w cm, minimum height=1.0cm, text width=\w cm-8pt, font=\footnotesize]
    at ($(c\k)+(\w/2,2.68)$) {Cohesion \t};}
\node[box=steellight, minimum width=2*\w cm+\g cm, minimum height=1.0cm, font=\footnotesize]
  at ($(c3)!0.5!(c4)+(\w/2,2.68)$) {Cohesion increased};

% ---------- row 3: the groups
\foreach \k/\n/\en/\de in {0/1/Primary shaping/Urformen,1/2/Forming/Umformen,2/3/Separating/Trennen,
                           3/4/Joining/Fügen,4/5/Coating/Beschichten,5/6/Properties/Stoffeigenschaft\\[-5pt]ändern}{
  \ifnum\k=0
    \node[boxstrong=amaccent, minimum width=\w cm, minimum height=1.7cm, text width=\w cm-8pt]
      (g\k) at ($(c\k)+(\w/2,1.2)$) {\n\quad\en\\[-1pt]{\scriptsize\mdseries\itshape\de}};
  \else
    \node[boxstrong=conv, minimum width=\w cm, minimum height=1.7cm, text width=\w cm-8pt]
      (g\k) at ($(c\k)+(\w/2,1.2)$) {\n\quad\en\\[-1pt]{\scriptsize\mdseries\itshape\de}};
  \fi}

% ---------- row 4: pictograms (each in a 2.4 x 1.6 cell, origin = bottom-left of cell)
\begin{scope}[shift={($(c0)+(0,-1.55)$)}]  % casting
  \draw[mould] (0.55,0.05) rectangle (1.85,0.85);
  \fill[white] (0.8,0.3) rectangle (1.6,0.85);
  \draw[steel, line width=0.8pt] (0.8,0.85) -- (0.8,0.3) -- (1.6,0.3) -- (1.6,0.85);
  \fill[mat=liquid] (0.8,0.3) rectangle (1.6,0.62);
  \draw[machine] (1.05,1.45) -- (1.55,1.45) -- (1.48,1.1) -- (1.12,1.1) -- cycle;
  \draw[liquid, line width=2pt] (1.2,1.1) .. controls (1.18,0.9) .. (1.2,0.62);
\end{scope}
\begin{scope}[shift={($(c1)+(0,-1.55)$)}]  % forging
  \draw[machinedark] (0.6,1.15) rectangle (1.8,1.45);
  \draw[machinedark] (0.6,0.05) rectangle (1.8,0.35);
  \fill[matsolid=plastic] (0.75,0.35) rectangle (1.65,1.15);
  \draw[flow=steel] (1.2,1.72) -- (1.2,1.5);
  \draw[motion] (0.6,0.75) -- (0.35,0.75); \draw[motion] (1.8,0.75) -- (2.05,0.75);
\end{scope}
\begin{scope}[shift={($(c2)+(0,-1.55)$)}]  % cutting
  \fill[matsolid=powder] (0.4,0.1) -- (2.0,0.1) -- (2.0,0.65) -- (1.3,0.65) -- (1.3,0.8) -- (0.4,0.8) -- cycle;
  \draw[machine] (1.15,0.85) rectangle (1.45,1.6);
  \draw[steel, line width=0.6pt] (1.15,1.05) -- (1.45,1.25) (1.15,1.3) -- (1.45,1.5);
  \draw[powder!70!black, line width=1pt] (1.55,0.75) .. controls (1.75,0.95) and (1.9,0.85) .. (1.85,1.1);
\end{scope}
\begin{scope}[shift={($(c3)+(0,-1.55)$)}]  % welding
  \fill[matsolid=powder] (0.3,0.2) rectangle (1.15,0.75);
  \fill[matsolid=powder] (1.25,0.2) rectangle (2.1,0.75);
  \fill[heat, draw=heat!60!black, line width=0.4pt] (1.0,0.75) .. controls (1.1,0.95) and (1.3,0.95) .. (1.4,0.75) -- cycle;
  \fill[heat] (1.15,0.2) rectangle (1.25,0.75);
  \draw[machine] (1.1,1.65) -- (1.3,1.65) -- (1.24,1.05) -- (1.16,1.05) -- cycle;
  \draw[energy, line width=0.8pt, decorate, decoration={zigzag,segment length=2.5pt,amplitude=1pt}] (1.2,1.05) -- (1.2,0.9);
\end{scope}
\begin{scope}[shift={($(c4)+(0,-1.55)$)}]  % coating
  \fill[matsolid=powder] (0.35,0.1) rectangle (2.05,0.65);
  \fill[ionised!70] (0.35,0.65) rectangle (2.05,0.8);
  \foreach \xx in {0.7,1.0,1.3,1.6} \draw[motion] (\xx,1.55) -- (\xx+0.08,0.88);
  \draw[machine] (1.0,1.6) rectangle (1.6,1.75);
\end{scope}
\begin{scope}[shift={($(c5)+(0,-1.55)$)}]  % heat treatment
  \fill[matsolid=powder] (0.55,0.35) rectangle (1.85,0.95);
  \foreach \xx in {0.75,1.2,1.65} \draw[heatwave] (\xx,0.0) -- (\xx,0.3);
  \foreach \xx in {0.75,1.2,1.65} \draw[heatwave] (\xx,1.5) -- (\xx,1.0);
\end{scope}

% ---------- row 5: conventional examples
\foreach \k/\t in {0/{casting\\injection moulding\\pressing + sintering\\electroforming},
                   1/{forging\\rolling\\deep drawing\\bending},
                   2/{turning, milling\\drilling, grinding\\laser cutting\\EDM},
                   3/{welding\\brazing\\adhesive bonding\\riveting},
                   4/{painting\\thermal spraying\\PVD, CVD\\electroplating},
                   5/{heat treatment\\hardening\\HIP\\shot peening}}{
  \node[lblc, anchor=north, text width=\w cm, font=\footnotesize, text=inksoft]
    at ($(c\k)+(\w/2,-1.65)$) {\t};}

% ---------- row 6: link with AM
\foreach \k/\t in {0/{\textbf{every AM process}:\\primary shaping,\\selective, layer\\by layer, no mould},
                   1/{bonding by plastic\\deformation: cold spray,\\friction stir, ultrasonic},
                   2/{support removal,\\finish machining,\\hybrid DED + milling},
                   3/{DED and wire-arc AM\\grew out of welding;\\fewer joints},
                   4/{cladding and cold spray:\\coatings grown into\\parts; repair},
                   5/{post-processing:\\stress relief, HIP,\\ageing, peening}}{
  \ifnum\k=0
    \node[box=amaccent, anchor=north, minimum width=\w cm, minimum height=2.05cm, text width=\w cm-8pt,
      font=\footnotesize, draw=amaccent, line width=1.1pt] at ($(c\k)+(\w/2,-3.45)$) {\t};
  \else
    \node[box=amaccent, anchor=north, minimum width=\w cm, minimum height=2.05cm, text width=\w cm-8pt,
      font=\footnotesize] at ($(c\k)+(\w/2,-3.45)$) {\t};
  \fi}
\node[font=\bfseries\footnotesize, text=amaccent, rotate=90] at (-0.35,-4.48) {link with AM};
\node[font=\bfseries\footnotesize, text=conv, rotate=90] at (-0.35,-1.6) {conventional};
\end{tikzpicture}
\end{document}
